% Generated draft: requires mathematical and source review.
\documentclass[11pt]{article}
\usepackage[margin=25mm]{geometry}
\usepackage{fontspec}
\setmainfont{Latin Modern Roman}
\usepackage{amsmath,amssymb,mathtools,unicode-math}
\setmathfont{Latin Modern Math}
\usepackage{enumitem,needspace,xurl,hyperref}
\setlength{\parindent}{0pt}
\setlength{\parskip}{5pt}
\setlength{\emergencystretch}{3em}
\newcommand{\N}{\mathbb{N}}
\newcommand{\Z}{\mathbb{Z}}
\newcommand{\Q}{\mathbb{Q}}
\newcommand{\R}{\mathbb{R}}
\newcommand{\C}{\mathbb{C}}
\newcommand{\F}{\mathbb{F}}
\newcommand{\A}{\mathbb{A}}
\newcommand{\PP}{\mathbb{P}}
\newcommand{\E}{\mathbb{E}}
\newcommand{\eps}{\varepsilon}
\newcommand{\dd}{\,\mathrm{d}}
\begin{document}
\section*{UnsolvedMath ID: 30002550}
\textbf{Source problem code:} catalogue entry
\subsection{Definitions and mathematical statement}
The supplied problem asks: ``Determine the minimax-optimal rates of convergence for estimation in sparse, high-dimensional varying-coefficient models.''

A minimax formulation requires a statistical experiment \(\{P_{\theta,n}:\theta\in\Theta_n\}\), where \(n\) denotes sample size, a parameter class \(\Theta_n\), an estimand \(T_n(\theta)\), and a nonnegative loss \(L_n\). Once these are specified, define the minimax risk by
\[
R_n(\Theta_n,T_n,L_n)
=\inf_{\widehat T_n}\sup_{\theta\in\Theta_n}
\mathbb E_{\theta,n}\!\left[L_n\bigl(\widehat T_n,T_n(\theta)\bigr)\right],
\]
where the infimum is over all measurable estimators based on the observations and taking values in the specified estimation space. A minimax-optimal rate is a positive sequence \(r_n\) for which there are constants \(0<c\leq C<\infty\), independent of \(n\) and of any dimension, sparsity, or regularity parameters over which uniformity is required, such that
\[
c\,r_n\leq R_n(\Theta_n,T_n,L_n)\leq C\,r_n
\]
throughout the specified asymptotic regime, for all sufficiently large \(n\). Any permitted dependence of the constants on fixed model parameters must also be specified.

\textbf{Requested determination.} Specify the intended sparse, high-dimensional varying-coefficient experiment and parameter class, and determine the corresponding minimax-optimal rate, including its dependence on sample size, dimension, sparsity, and coefficient regularity where applicable. The display above is a general framework for interpreting the request, not a statistical model supplied by the source record.

\textbf{Unresolved qualifications.} The record does not specify the observation equation, sampling and noise assumptions, covariate design, coefficient-function domain or boundary conditions, smoothness class, meaning of sparsity, estimation target, loss, or growth regime. Consequently it does not determine a unique minimax risk or rate. Its references to homogeneous cases and to a general exact minimax-rate question are also not accompanied by definitions of those cases or their hypotheses.
\subsection{Short English statement}
Determine the best possible estimation rates for sparse, high-dimensional varying-coefficient models. A precise rate requires an explicit statistical model, parameter class, estimation target, and loss; these are missing from the supplied record. Its supplied partially solved status and dated literature assessment require independent mathematical review.
\subsection{Sources}
\begin{itemize}
\item Supplied local registry record: \texttt{data/unsolvedmath/problems.json}, permanent UnsolvedMath ID \(30002550\), source problem code \texttt{catalogue entry}, title ``Minimax Rates for Sparse Varying-Coefficient Models.'' The record belongs to the Oberwolfach Reports Open Problems collection and lists proposed year 2014; no individual proposer is supplied.
\item Supplied workshop provenance: \emph{Adaptive Statistical Inference} (2014), Oberwolfach Reports, pp.~721--779, DOI 10.4171/owr/2014/13. \url{https://doi.org/10.4171/owr/2014/13}. The record describes these reports as documenting mathematical workshops and their open-problem discussions; no problem-specific page or named contributor is given.
\item Supplied literature reference: Klopp--Pensky (2015). \url{https://doi.org/10.1214/15-AOS1309}.
\item Supplied literature reference: Klopp--Pensky preprint. \url{https://arxiv.org/abs/1312.4087}.
\item Supplied literature assessment, checked on \texttt{2026-08-21}: the record labels the problem partially solved and states that Klopp and Pensky established minimax lower bounds and adaptive estimators attaining them up to constants in homogeneous cases or a logarithmic factor otherwise, leaving the general exact minimax-rate question not fully settled. This is supplied dated curation derived from the OWR collection, not an independently verified assessment. Its time-sensitive status and cited claims require independent verification.
\end{itemize}
\end{document}
